\documentclass[12pt]{article}
\usepackage[utf8]{inputenc}
\usepackage[T1]{fontenc}
\usepackage{hyperref}
\title{Mpalitjanh Language Resources for K-12 Curriculum Development}
\author{Compiled by Claude Sonnet 5.5 (claude-sonnet-5-5) for Andr\'as Kornai}
\date{}
\begin{document}
\maketitle

\section{Introduction}
Mpalitjanh (also spelled Mpalityan, Mbalidjan) is an extinct Pama-Nyungan language of the Paman group on the Cape York Peninsula, Queensland. Sources I saw treat it as a dialect of the same language as Luthigh, and associate it with a sub-group of the Winduwinda people. I found no date of extinction. Very little material exists, and I found no teaching resources.

\section{Primary Educational Resources}
I found no primers, textbooks or school programmes.

\section{Linguistic and Reference Materials}
\begin{itemize}
\item Wikipedia, Mpalitjanh dialect: \url{https://en.wikipedia.org/wiki/Mpalitjanh_dialect}
\item Wikipedia, Luthigh language (the related language of which Mpalitjanh is treated as a dialect): \url{https://en.wikipedia.org/wiki/Luthigh_language}
\item Hale, Kenneth L. 1976, on Northern Paman phonology (Glottolog bibliographic record): \url{https://glottolog.org/resource/reference/id/306792}
\item Glottolog bibliographic record (ozbib:2457): \url{https://glottolog.org/resource/reference/id/ozbib:2457}
\item Luthigh entry in the Handbook of Australian cultures (PDF): \url{https://www.hgcosmos.org/assets/culture/AZ/L/Luthigh.pdf}
\end{itemize}

\section{Religious and Community Resources}
I found no religious texts or dedicated community organisations. One community profile page mentions the language: \url{https://aboutthenorth.au/styled-59/styled-184/styled-238/styled-446/}

\section{Digital and Audio Resources}
Wikitongues has a language page for xpj. I did not see any audio or video on it in the search results: \url{https://wikitongues.org/languages/xpj/}

\section{Curriculum Development Context}
Documentation is thin, so any curriculum would be a heritage-awareness programme built with the descendant community (the Winduwinda people and neighbouring Cape York communities). Any wordlists should be sought through AIATSIS (code Y25) and the community, and used only with community consent. No vocabulary is given here.

\end{document}